\subsection{Sonic Pi}\label{subsec:sonic-pi}

\textbf{Sonic Pi} began as a 3-month pilot project in \textit{November 2013}, funded by the \textit{Broadcom
Foundation}.
It aimed to teach school students with no prior programming experience core computational concepts through a live coding
environment, as well as enable music composers to deliver live coded performances~\cite{Aaron_Blackwell_2013,
Aaron_2016}.

This language was developed to run on the \textit{Raspberry Pi}, having to work on tight hardware constraints, where
time and memory overhead needed to be managed carefully.
It's core architecture was designed around an earlier project called \textit{Overtone}, but which had too steep of a
learning curve for music performers and had massive performance issues on the \textit{Raspberry Pi}~\cite{Aaron_2016}.

%% ---------------------------------------------------------------------------------------------------------------------

\subsubsection{Key Features}

To address the tight constraints of having to be optimized to run on low end hardware such as the \textit{Raspberry Pi},
the language was designed and embedded in \textit{Ruby}.

The system is designed in terms of the \textit{SuperCollider} sound synthesis engine, allowing the control
and manipulation of synthesizers in real time.
\textit{SuperCollider} acts as the core music synthesis backend for \textbf{Sonic Pi}, as it has no synthesis
capabilities of its own~\cite{Aaron_Blackwell_2013}.

\textbf{Sonic Pi} follows a \textit{sequential execution model}, where musical events are generated in the order in
which instructions are executed.
We can see this in action in the following code snippets explaining its core methods.

The simplest way to play a sound is through the \texttt{play} method, which takes a \textit{MIDI} note as its argument.
Notes can be separated and spaced out between each other using the \texttt{sleep} method, which takes a floating point
number as its argument, representing the time in seconds to wait until playing the next sound.
Notes can be combined into patterns by using simple lists, and they can be played sequentially using the
\texttt{play\_pattern} method~\cite{Aaron_Blackwell_2013}.

Example snippet:
\begin{lstlisting}[label={lst:sonic-pi-1}]
  play 60;
  sleep 0.5;
  play_pattern [50, 55, 60];
\end{lstlisting}

\textbf{Sonic Pi} also introduces simple control flow structures such as \textit{loops} and
\textit{conditional statements}:
\begin{lstlisting}[label={lst:loops-and-conditional-statements}]
  10.times do
    play 65;
    sleep 0.25;
    ...
  end

  if num < 0.5
    play 60
  else
    play 70
  end
\end{lstlisting}

One more powerful feature is the ability to stack options or filters on a certain node, by specifying the name of the
option and a value, enabling complex sound modulation and more granularity and control over notes.
Examples of these options include \texttt{attack}, \texttt{release}, \texttt{sustain} and many more.

This simplistic model enables an easier grasp over musical compositional concepts for users with non-technical
backgrounds, and is the main reason why it had such a positive impact in teaching sessions in middle schools in the
UK~\cite{Aaron_2016, Aaron_Blackwell_2013}.

%% ---------------------------------------------------------------------------------------------------------------------

\subsubsection{Limitations}

\textbf{Sonic Pi} functions more as an \textit{OSC} client that communicates to a subset of the \textit{SuperCollider}
API\@.
This subset, although small, is just enough for triggering sounds, controlling synthesizers and clearing data between
runs.
This, however, leaves \textbf{Sonic Pi} incapable of synthesizing sound on its own, and thus relies on an external
synthesis engine~\cite{Aaron_Blackwell_2013}.

Due to it being primarily and educational tool and its sequential execution model, it offers limited compositional
scalability and structural abstraction.
Complex layering capabilities are sacrificed for better ease of use and mental models, leaving the language weaker than
\textit{Overtone}, on which it was based on.

%% ---------------------------------------------------------------------------------------------------------------------

\subsubsection{Relevance to Our DSL}

The syntax of this language is both clean and concise, thus making it easier to map abstract concepts from musical
theory into textual form.
It creates the building blocks necessary to construct more complex structures, providing a good basis for our DSL's key
syntactical model.
However, the sequential execution model can make it more difficult to explicitly represent complex temporal
relationships between independent musical structures, which motivates the timeline-based approach proposed in this
thesis.
